\title{xLSTM: Extended Long Short-Term Memory}

\begin{abstract}
The foundational mechanisms behind Long Short-Term Memory (LSTM) networks—specifically, the constant error carousel and gating—were established in the 1990s. Over subsequent decades, LSTMs have proven highly robust, serving as the backbone for significant milestones in deep learning, notably as precursors to the earliest Large Language Models (LLMs). The introduction of Transformer architectures, featuring inherently parallel self-attention mechanisms, however, initiated a transformative phase in sequence modeling that quickly surpassed LSTMs when scaled. This leads us to a pertinent question: what levels of language modeling performance can be achieved if LSTMs are scaled to billions of parameters, integrating contemporary methods from LLM research while addressing traditional shortcomings? To this end, we first propose an exponential gating mechanism supplemented by suitable normalization and stabilization strategies. We then redesign the LSTM’s memory operations to develop: (i) sLSTM, which employs scalar memory and updates alongside a novel memory mixing strategy; and (ii) mLSTM, which is structured for full parallelism through matrix memory and an updated covariance rule. By embedding these enhanced LSTM variants within residual network structures, we obtain xLSTM blocks, which are then assembled into complete xLSTM models. The combination of exponential gating and restructured memory notably enhances the xLSTM, enabling it to achieve competitive or superior results compared to advanced Transformers and State Space Models in both effectiveness and scalability.\\
Code available at: \url{https://github.com/NX-AI/xlstm}
\end{abstract}

\section{Introduction}
\label{sec:intro}

Concepts underlying the Long Short-Term Memory (LSTM) architecture~\citep{Hochreiter:91,Hochreiter:97nips,Hochreiter:97}, namely the constant error carousel and gating mechanisms, were proposed as solutions to address the vanishing gradient issue observed in recurrent neural networks~\citep{Hochreiter:91,Hochreiter:00book}:

\begin{align}
\label{eq:lstm_idea}
  c_t \ = \ f_t \ c_{t-1} \ + \ 
          i_t \  z_t \ , \quad  
          h_t \ = \ o_t \ \psi({c_t}) \ . 
\end{align}

The constant error carousel functions by additively modifying the cell state $c_{t-1}$ (depicted in green) through cell inputs $z_t$, with the magnitude of this update being regulated by sigmoid gates (shown in blue). The input gate $i_t$ and forget gate $f_t$ are responsible for governing this process, whereas the output gate $o_t$ determines the emission from the memory cell, specifically the hidden state $h_t$. After normalization or squashing of the cell state using $\psi$, the output gate subsequently modulates this processed state to yield the hidden state.

LSTMs have found widespread application across multiple fields \citep{Hochreiter:01,Hochreiter:07,Schmidhuber:15}, maintaining dominance in text generation tasks until the introduction of Transformers in 2017~\citep{Vaswani:17}. Their proficiency has been validated in a range of sequence-driven problems, including but not limited to text production~\citep{Graves:13,Karpathy:15blog}, handwriting synthesis~\citep{Graves:13}, sequence-to-sequence language translation~\citep{Sutskever:14nips}, program analysis~\citep{Zaremba:14arxiv}, image caption creation~\citep{Karpathy:15,Hossain:19}, code generation~\citep{Karpathy:15blog}, rainfall-runoff forecasting~\citep{Kratzert:18,Kratzert:19}, as well as hydrological modeling for flood prediction~\citep{Nearing:24}. In the sphere of reinforcement learning, LSTMs stand out as the leading sequence modeling framework, exemplified by their use in AlphaStar for StarCraft II~\citep{Vinyals:17short}, OpenAI Five for Dota 2~\citep{OpenAI:19}, and nuclear fusion magnetic controller models~\citep{Degrave:22short}. Their strength lies in their ability to internalize complex abstractions, efficiently capturing and retaining semantic content within their memory structures~\citep{Karpathy:15blog}; this capacity is illustrated by the discovery of neurons dedicated to numeracy and syntactic processing~\citep{Lakretz:19}, linguistic comprehension~\citep{Bau:19}, and sentiment analysis~\citep{Radford:17arxiv}. Even today, LSTMs continue to underpin critical, impactful systems~\citep{Degrave:22short,Nearing:24}, underscoring their enduring relevance and robustness.

Although LSTMs have achieved remarkable results, they face three primary drawbacks: (i) Difficulty in revising previously made storage decisions. This shortcoming is illustrated with the {\it Nearest Neighbor Search} task (refer also to Appendix~\ref{sec:appNearestNeighborSearch}). Given a reference vector, the sequence must be traversed to identify the most similar vector and return its corresponding value at the end of the sequence. The left side of Figure~\ref{fig:lstmProblems} presents the mean squared error for this task. LSTMs exhibit challenges in updating a stored value when encountering a more relevant match later in the sequence, whereas our proposed xLSTM overcomes this issue by implementing exponential gating. (ii) Restricted capacity for information storage, as information must be encoded into single-dimensional cell states. This restriction is highlighted through the {\it Rare Token Prediction} problem. The right side of Figure~\ref{fig:lstmProblems} displays token prediction perplexities on Wikitext-103~\citep{Merity:17} according to token frequency partitions. Due to their constrained storage, LSTMs struggle more with infrequent tokens, a challenge addressed by xLSTM’s matrix-based memory. (iii) Inherent difficulties with parallel computation because of the interdependence between hidden states across time steps, resulting from hidden-to-hidden connections that necessitate strictly sequential computation.

The challenges inherent in LSTM architectures have catalyzed the rise of Transformers~\citep{Vaswani:17} within the domain of language modeling. This raises an important question: if we address these limitations and increase the scale of LSTMs to that of modern Large Language Models, what levels of performance might we observe in language modeling tasks?

\section{Extended Long Short-Term Memory}
\label{sec:xLSTM}

In addressing the constraints inherent to LSTM architectures, Extended Long Short-Term Memory (xLSTM) makes two principal changes to the formulation in Equation~\eqref{eq:lstm_idea}. Specifically, xLSTM incorporates exponential gating and develops new memory architectures, resulting in two novel variants: (i) sLSTM (see Section~\ref{sec:sLSTM}), which uses a scalar memory, applies scalar updates, and implements memory mixing; and (ii) mLSTM (see Section~\ref{sec:mLSTM}), which maintains a matrix-based memory structure and utilizes an outer product update scheme designed for covariance, allowing for full parallelization. Exponential gating serves to bolster the effectiveness of both the sLSTM and mLSTM models. In pursuit of parallel computation, mLSTM omits the memory mixing step, thus forgoing hidden-to-hidden recurrent connections. Both architectures, sLSTM and mLSTM, permit generalization to settings involving several memory cells; within sLSTM, this enables memory mixing among cells. Furthermore, sLSTM can be configured with multiple heads that do not interact via memory mixing between heads, but allow memory mixing among cells within each head. The combination of multiple heads with exponential gating in sLSTM yields an alternative paradigm for memory mixing. For mLSTM, the concepts of multiple heads and multiple cells coincide.

By incorporating these novel LSTM variants within residual block modules, we obtain xLSTM blocks (refer to Section~\ref{sec:xLSTMblock}). Building architectures by stacking these xLSTM blocks in a residual manner yields xLSTM architectures (see Section~\ref{sec:xLSTMarchitecure}). Figure~\ref{fig:xlstm_sketch} illustrates the overall xLSTM architecture and its constituent components.

\subsection{Review of the Long Short-Term Memory}
\label{sec:orgLSTM}

The foundational concept of LSTM~\citep{Hochreiter:91,Hochreiter:97nips,Hochreiter:97} established the scalar memory cell as a pivotal unit for both computation and information retention, designed specifically to mitigate the vanishing gradient problem~\citep{Hochreiter:91,Hochreiter:00book} by leveraging the mechanism known as the constant error carousel, which is realized in the cell state update. The memory cell structure incorporates three gating mechanisms: input gate, output gate, and forget gate—the last of which was introduced by \citet{Gers:00}. 
The LSTM memory cell is updated at each time step $t$ according to the following rules:

\begin{align}
c_t \ &= \  f_t \ c_{t-1} \ + \ i_t \ z_t &  & &\text{cell state} \\
h_t  \ &= \ o_t \ \tilde{h}_t \ , & \tilde{h}_t \ &= \ \psi \left( c_t\right)
  &\text{hidden state} \\
z_t \ &= \ \varphi \left( \tilde{z}_t \right) \ , 
  &\tilde{z}_t \ &=  \ \mathbf{w}^\top_{z} \ \mathbf{x}_t \ + \
  r_{z}  h_{t-1} \ + \  b_{z} \ \
  &\text{cell input} \\
i_t \ &= \ \sigma \left( \tilde{i}_t  \right) \ , 
  &\tilde{i}_t \ &= \ \mathbf{w}^\top_{i} \ \mathbf{x}_t \ + \
  r_{i}  \ h_{t-1} \ + \  b_{i} \ \
  &\text{input gate} \\
f_t \ &= \ \sigma \left(  \tilde{f}_t \right) \ , 
  &\tilde{f}_t \ &= \ \mathbf{w}^\top_{f} \ \mathbf{x}_t  \ + \
  r_{f}  \ h_{t-1} \ + \  b_{f} \ \
  &\text{forget gate} \\
o_t \ &= \ \sigma \left( \tilde{o}_t \right) \ , 
  &\tilde{o}_t  \ &= \ \mathbf{w}^\top_{o} \ \mathbf{x}_t \ + \
  r_{o}  \ h_{t-1} \ + \  b_{o} \ \
  &\text{output gate} 
\end{align}

The vectors $\mathbf{w}_{z}$, $\mathbf{w}_{i}$, $\mathbf{w}_{f}$, and $\mathbf{w}_{o}$ represent the input weights that connect the input $\mathbf{x}_t$ to the cell input, as well as the input, forget, and output gates, respectively. The recurrent weights $r_{z}$, $r_{i}$, $r_{f}$, and $r_{o}$ denote the connections from the previous hidden state $h_{t-1}$ to the cell input and the various gates in the same order. The associated bias terms are given by $b_{z}$, $b_{i}$, $b_{f}$, and $b_{o}$. The functions $\varphi$ and $\psi$ act as the activation functions for the cell input and the hidden state, typically instantiated as $\tanh$. Notably, $\psi$ serves to constrain or normalize the cell state, preventing it from becoming unbounded. Each gate uses a sigmoid activation, $\sigma \left( x \right)= 1/(1+\exp(-x))$. Subsequent developments introduced the aggregation of multiple scalar memory cells $\mathbf{c}_t \in \mathbb{R}^{d}$ into a single vector, which enables the application of recurrent weight matrices $\mathbf{R} \in \mathbb{R}^{d \times d}$ to intermix the outputs from different memory cells~\citep{Greff:15}. Further information is available in Appendix~\ref{sec:apporgLSTM}. Ablation experiments have indicated that each element of the memory cell architecture plays an essential role~\citep{Greff:15}.

\subsection{sLSTM}
\label{sec:sLSTM}

To enable LSTMs to modify storage choices dynamically, we propose the incorporation of exponential gates (depicted in red), along with mechanisms for normalization and stabilization. Specifically, the input and forget gates may employ exponential activation functions. For normalization purposes, we define a normalizer state that aggregates the products of the input gate and the sequence of subsequent forget gates.

The scalar sLSTM forward computation is as follows:

\begin{align}
\label{eq:slstmforward}
c_t \ &= \  f_t \ c_{t-1} \ + \ i_t \ z_t & & 
 &\text{cell state} \\
n_t \ &= \  f_t \ n_{t-1} \ + \ i_t  & & 
 &\text{normalizer state} \\
h_t \ &= \ o_t \ \tilde{h}_t \ ,
  &\tilde{h}_t \ &= \ c_t / n_t
&\text{hidden state} \\
z_t \ &= \ \varphi \left( \tilde{z}_t \right) \ , 
  &\tilde{z}_t \ &=  \ \mathbf{w}^\top_{z} \ \mathbf{x}_t \ + \
  r_{z}  h_{t-1} \ + \  b_{z}
  &\text{cell input} \\
i_t \ &= \ \!\exp \left( \tilde{i}_t  \right) \ , 
  &\tilde{i}_t \ &= \ \mathbf{w}^\top_{i} \ \mathbf{x}_t \ + \
  r_{i}  \ h_{t-1} \ + \  b_{i}
  &\text{input gate} \\
f_t \ &= \ \sigma \left(  \tilde{f}_t \right) \ \text{OR} \
\!\exp \left(  \tilde{f}_t \right) \ ,
  &\tilde{f}_t \ &= \ \mathbf{w}^\top_{f} \ \mathbf{x}_t  \ + \
  r_{f}  \ h_{t-1} \ + \  b_{f}
  &\text{forget gate} \\
o_t \ &= \ \sigma \left( \tilde{o}_t \right) \ , 
  &\tilde{o}_t  \ &= \ \mathbf{w}^\top_{o} \ \mathbf{x}_t \ + \
  r_{o}  \ h_{t-1} \ + \  b_{o}
  &\text{output gate} 
\end{align}

We adapt the foundational LSTM gating strategies—including gating mechanisms that depend on either the input and/or hidden states, accompanied by a bias term—to these novel model structures. Since exponential activation functions have the potential to generate excessively large outputs, potentially resulting in numerical overflows, we introduce an extra state variable, \mbox{$m_t$~\citep{Milakov:18arxiv}}, to help regulate and stabilize the gate activations:

\begin{align}
\label{eq:slstmstabil}
m_t \ &= \ \max \left( \log ( f_t ) + m_{t-1} , \log ( i_t ) \right) &\text{stabilizer state} \\
i'_t \ &= \ \!\exp \left( \log \left ( i_t \right) - m_t \right) \ = \ \!\exp \left( \tilde{i}_t - m_t \right) \ \, 
  &\text{stabil. input gate} \\
f'_t \ &= \ \!\exp \left( \log \left( f_t \right) + m_{t-1} - m_t \right) \ \, 
  &\text{stabil. forget gate}
\end{align}

As demonstrated in Appendix~\ref{sec:appsLSTM}, substituting $f_t$ with $f'_t$ and $i_t$ with $i'_t$ during the forward computation does not alter either the overall network output or the gradients of the loss with respect to its parameters.

\paragraph{New Memory Mixing.} Similar to traditional LSTM architectures (see Appendix~\ref{sec:appsLSTM}), sLSTM is capable of supporting multiple memory cells. The use of several memory cells facilitates memory mixing through recurrent connections $\mathbf{R}_{\mathbf{z}}$, $\mathbf{R}_{\mathbf{i}}$, $\mathbf{R}_{\mathbf{f}}$, $\mathbf{R}_{\mathbf{o}}$, which transmit information from the hidden state vector $\mathbf{h}$ to both the memory cell input $\mathbf{z}$ and the corresponding gates $\mathbf{i}$, $\mathbf{f}$, and $\mathbf{o}$. A distinguishing characteristic of this memory mixing approach is the incorporation of exponential gating. In the revised sLSTM model, each head may perform memory mixing independently, whereas such mixing does not occur between different heads. By integrating heads in sLSTM along with exponential gating, a novel framework for memory mixing is achieved.

\subsection{mLSTM}
\label{sec:mLSTM}

In order to improve the storage capability of LSTMs, we transition from representing the memory cell as a scalar $c \in \mathbb{R}$ to a matrix form $\mathbf{C} \in \mathbb{R}^{d \times d}$. This modification means that retrieval operations are carried out through matrix multiplication. At each time step $t$, the network is designed to store two vectors: a key $\mathbf{k}_t \in \mathbb{R}^d$ and an associated value $\mathbf{v}_t \in \mathbb{R}^d$, following nomenclature from the Transformer architecture. Subsequently, at a future time $t+\tau$, we seek to recover the value $\mathbf{v}_t$ by applying a query vector $\mathbf{q}_{t+\tau} \in \mathbb{R}^d$. This configuration corresponds to the framework used in Bidirectional Associative Memories (BAMs) \citep{Kohonen:72,Anderson:72,Nakano:72,Anderson:77}.

To encode the key-value association, we utilize the covariance update mechanism described in~\citep{Sejnowski:77,Dayan:91}:

\begin{align}
\mathbf{C}_t \ &= \ \mathbf{C}_{t-1} \ + \ \mathbf{v}_{t} \ \mathbf{k}_t^\top \ .
\end{align}

We consider a scenario in which inputs are first normalized via layer-norm prior to being mapped onto key and value representations, resulting in vectors with zero mean. According to~\citep{Dayan:91}, the covariance update mechanism achieves optimality by maximizing the separability among retrieved binary vectors; this, in turn, maximizes the signal-to-noise ratio. Permitting only pairwise retrieval interactions, similar to attention mechanisms, and accepting their quadratic complexity enables even greater separability~\citep{Krotov:16,Krotov:17,Ramsauer:21}. This form of covariance updating aligns with the methodology used in Fast Weight Programmers~\citep{Schmidhuber:92ncfastweights,Schlag:21}, which were subsequently augmented to include a constant decay factor applied to $\mathbf{C}_{t-1}$ and a constant learning rate governing updates of the form 
$\mathbf{v}_{t} \mathbf{k}_t ^\top$~\citep{Ba:16}.
Building on this idea, we embed the covariance update rule within the LSTM structure: the forget gate acts as the decay rate, the input gate serves as the learning rate, and the output gate modulates the magnitude of the retrieved vector.

In the context of this matrix memory, the normalizer state is constructed as a weighted aggregation of key vectors, with each vector being modulated by the input gate along with all subsequent forget gates. This configuration allows the normalizer state to track the influence or magnitude of the various gates over time. Given that the dot product between the query and the normalizer state might approach zero, we take the absolute value of this product and ensure it does not fall below a specified threshold (commonly set to 1.0), following the methodology described in~\citep{Sun:23arxiv}. The mLSTM computes the forward pass as follows:

\begin{align}
\label{eq:mlstm_recurrent_begin}
\mathbf{C}_t \ &= \  f_t \ \mathbf{C}_{t-1} \ + \ 
  i_t \ \mathbf{v}_t \ \mathbf{k}_t^\top & &  &\text{cell state} \\
\mathbf{n}_t \ &= \  f_t \ \mathbf{n}_{t-1} \ + \ 
  i_t \ \mathbf{k}_t & &  &\text{normalizer state} \\
\mathbf{h}_t  \ &= \ \mathbf{o}_t \ \odot \ \tilde{\mathbf{h}}_t
\ , \qquad \qquad 
\tilde{\mathbf{h}}_t \ = \   \mathbf{C}_t \mathbf{q}_t \ / \ 
  \max \left\{ |\mathbf{n}_t^\top \mathbf{q}_t|, 1 \right\} 
   &&&\text{hidden state} \\
\mathbf{q}_t \ &= \ \mathbf{W}_q \ \mathbf{x}_t \ + \ \mathbf{b}_q  & &  &\text{query input} \\
\mathbf{k}_t \ &= \ \frac{1}{\sqrt{d}} \mathbf{W}_k \ \mathbf{x}_t \ + \ \mathbf{b}_k  & &  &\text{key input} \\
\mathbf{v}_t \ &= \ \mathbf{W}_v \ \mathbf{x}_t \ + \ \mathbf{b}_v  & &  &\text{value input} \\
i_t \ &= \ \!\exp \!  \! \left( \tilde{i}_t  \right) \ , 
 \qquad \qquad \ \ \ \ \,
  \tilde{i}_t \ = \ \mathbf{w}^\top_{i} \ \mathbf{x}_t \ + \  b_{i} \ \ \ 
  &&&\text{input gate} \\
f_t \ &= \ \sigma \!  \left(  \tilde{f}_t \right) \ \text{OR} \ \!\exp \!  \! \left(  \tilde{f}_t \right) , \ \ \: \!
\tilde{f}_t \ = \ \mathbf{w}^\top_{f} \ \mathbf{x}_t  \ + \
  b_{f}
  &&& \text{forget gate} \\
\label{eq:mlstm_recurrent_end}
\mathbf{o}_t \ &= \ \sigma \left( \tilde{\mathbf{o}}_t \right) \ , \qquad \qquad \quad \ \ \ \,
  \tilde{\mathbf{o}}_t  \ = \ \mathbf{W}_{\mathbf{o}} \ \mathbf{x}_t \ + \
  \mathbf{b}_{\mathbf{o}}
  &&&\text{output gate} 
\end{align}

Similar to the original LSTM, mLSTM is capable of incorporating several memory cells. In the context of mLSTM, employing multiple heads is functionally the same as utilizing multiple cells, due to the absence of memory mixing. To ensure the stability of the exponential gates in mLSTM, we apply the stabilization strategies established for sLSTM (refer to Equation~\ref{eq:slstmstabil}). 
Because mLSTM does not involve mixing between memory states, its recurrence relation can be rewritten in a parallelizable form. Additional information can be found in Appendix~\ref{sec:appmLSTM}.

\subsection{xLSTM Architecture}
\label{sec:xLSTMarch}

\paragraph{xLSTM Blocks.}
\label{sec:xLSTMblock}
The core function of an xLSTM block is to perform a nonlinear summary of preceding sequence elements within a high-dimensional latent space, facilitating the disambiguation of distinct historical trajectories or contexts. Such differentiation among various histories is essential for accurately forecasting subsequent elements, such as predicting the following token in a sequence. To underpin this design choice, we invoke Cover’s Theorem~\citep{Cover:65}, which asserts that patterns which are nonlinearly embedded into higher-dimensional spaces tend to be more amenable to linear separation compared to those in lower-dimensional representations. We analyze two principal architectures for residual blocks: (i) A residual block employing post up-projection, as exemplified by Transformers, first compresses past information nonlinearly within the original dimensionality, then projects the result linearly into a higher-dimensional space, applies a nonlinear activation, and subsequently projects back linearly to the initial dimensionality; this architecture is visualized in the left portion of Figure~\ref{fig:Backbones} and in the third column of Figure~\ref{fig:xlstm_sketch}. A more granular depiction is provided in Figure~\ref{fig:appsLSTM_detailed} in the appendix. (ii) A residual block with pre up-projection, similar to the approach adopted in State Space Models, initially maps inputs linearly into a high-dimensional space,

The model first condenses historical information in a non-linear manner within a high-dimensional representation, after which it applies a linear transformation to return to the original feature space. When an xLSTM block incorporates an sLSTM component, the up-projection occurs after the main computation. In contrast, for xLSTM blocks with an mLSTM, the up-projection is performed beforehand, as the memory size increases in the high-dimensional context. Additional explanations are provided in the left panel of Figure~\ref{fig:Backbones} and the third column of Figure~\ref{fig:xlstm_sketch}, with further elaboration found in Figure~\ref{fig:appsLSTM_detailed} within the appendix.

\paragraph{xLSTM Architecture. }
\label{sec:xLSTMarchitecure}
The xLSTM architecture is formed by stacking modular components with residual connections~\citep{Srivastava:15,He:16}. Our approach adopts the widely utilized pre-LayerNorm~\citep{Ba:16b} residual backbone, aligning with the architectures of current Large Language Models. This structure is illustrated in the final column of Figure~\ref{fig:xlstm_sketch}.

\subsection{Memory and Speed Considerations}

Unlike Transformers, xLSTM networks exhibit linear computational complexity and maintain constant memory requirements regardless of sequence length. The compressive nature of xLSTM memory makes it particularly advantageous for deployment in industrial scenarios and on edge devices.

The mLSTM’s memory operates without the need for additional parameters; however, it incurs significant computational costs due to both its $d \times d$ memory matrix and the corresponding $d \times d$ update. This design choice reflects a balance between enhancing memory capacity and managing computational demands. Despite this, these operations can be efficiently parallelized on GPUs, which means they contribute minimally to the overall execution time.

Although mLSTM can be parallelized similarly to FlashAttention~\citep{Dao:22,Dao:23} or GLA~\citep{Yang:23arxiv}, sLSTM cannot be parallelized because of the memory mixing introduced by the hidden-hidden interactions. Nevertheless, we created an efficient CUDA-based implementation featuring GPU memory optimizations down to the register level, resulting in a runtime that is usually less than twice as slow as that of mLSTM.

\section{Related Work}
\label{sec:related}

\paragraph{Linear Attention.} Numerous approaches have been developed to address the Transformer's quadratic computational cost with respect to context length, aiming to achieve linear scaling in attention mechanisms. The Synthesizer model dispenses with explicit token--token interactions by directly learning synthetic attention weights~\citep{Tay:20arxiv}. Linformer approximates self-attention using a low-rank matrix representation and further exploits linear approximations~\citep{Wang:20arxiv}. The Linear Transformer reformulates the attention computation so that it becomes linear in context size \citep{Katharopoulos:20}. Performer utilizes positive orthogonal random features to provide a linear-time approximation of the softmax attention~\citep{Choromanski:21}. Moreover, attention has been supplanted by highly efficient long convolutional operations in models such as Structured Global Convolution (SGConv)~\citep{Li:22} and the Hyena Hierarchy~\citep{Poli:23}.

\paragraph{State Space Models.} In recent times, State Space Models (SSMs) have garnered significant attention due to their linear scalability with respect to sequence length and their strong empirical results, often rivaling those of Transformer-based architectures. Early influential works include the Structured State Space sequence model (S4)~\citep{Gu:21}. Subsequent developments in this area have produced architectures such as the Diagonal State Space (DSS) model~\citep{Gupta:22}, Gated State Space (GSS) models~\citep{Mehta:22}, the S5 model~\citep{Smith:22}, Bidirectional Gated SSM (BiGS)~\citep{Wang:22}, the H3 model~\citep{Fu:23}, and Mamba~\citep{Gu:24arxiv}.

\paragraph{Recurrent Neural Networks.} Recurrent Neural Networks (RNNs) have emerged as alternatives to Transformer-based architectures and attention mechanisms, particularly for their linear scalability with respect to sequence length. RNNs incorporating Deep Linear Recurrent Units (LRUs) have delivered notable advances in language modeling~\citep{Orvieto:23, De:24arxiv}. Similarly, models such as the Hierarchically Gated Linear RNN (HGRN)~\citep{Qin:23} and its successor HGRN2~\citep{Qin:24arxiv} have also demonstrated strong performance. Among RNN-based large language models, RWKV~\citep{Peng:23arxivshort,Peng:24arxivshort} stands out for achieving results comparable to those of Transformer models.

\paragraph{Gating.}
Gating, a central mechanism introduced in LSTM, has been independently rediscovered and further explored in a variety of recent methods. Variants of gating mechanisms have been incorporated into architectures such as HGRN~\citep{Qin:23}, HGRN2~\citep{Qin:24arxiv}, Gated Linear Attention (GLA)~\citep{Yang:23arxiv}, Gated State Space (GSS) models~\citep{Mehta:22}, Bidirectional Gated SSM (BiGS)~\citep{Wang:22}, Moving Average Equipped Gated Attention (MEGA)~\citep{Ma:22}, RWKV~\citep{Peng:23arxivshort}, and Mamba~\citep{Gu:24arxiv}.

\paragraph{Covariance Update Rule.} In order to improve storage capabilities, we integrated a matrix memory into the mLSTM cell using a covariance update rule. This update mechanism underlies various other architectures, including Fast Weight Programmers \citep{Schmidhuber:92ncfastweights,Schlag:21}, RWKV-5 and RWKV-6~\citep{Peng:24arxivshort}, Retention~\citep{Sun:23arxiv}, Linear Transformer~\citep{Katharopoulos:20}, and HGRN2~\citep{Qin:24arxiv}.

\paragraph{Most Related.} From a conceptual standpoint, the models that most closely resemble xLSTM are Retention~\citep{Sun:23arxiv}, RWKV~\citep{Peng:23arxivshort,Peng:24arxivshort}, and HGRN2~\citep{Qin:24arxiv}. These models employ mechanisms such as matrix-based memory and gating. Nevertheless, unlike the recently introduced sLSTM, these prior methods do not support memory mixing. The ability to mix memory is crucial for addressing state tracking challenges, making LSTMs inherently more powerful and expressive compared to State Space Models (SSMs) and Transformers~\citep{Merrill:24,Deletang:23}. Robust state tracking is essential for tasks like code evaluation or maintaining awareness of entities within extensive narratives.

\paragraph{Residually Stacking Architectures.} As with the majority of present-day large-scale deep learning systems, xLSTM models employ a residual stacking strategy when assembling their foundational modules~\citep{Srivastava:15,He:16}.

This design paradigm has facilitated breakthroughs in deep convolutional neural networks~\citep{He:16} as well as in the Transformer framework~\citep{Vaswani:17}. The Transformer has become the key architectural choice underlying Large Language Models (LLMs) such as GPT-3~\citep{Brown:20short}, ChatGPT~\citep{Schulman:22short}, GPT-4~\citep{Achiam:23gpt4short}, Megatron-LM~\citep{Shoeybi:19}, Gopher~\citep{Rae:21short}, ERNIE 3.0 Titan~\citep{Wang:21arxivshort}, GLaM~\citep{Du:21glamshort}, Chinese M6~\citep{Lin:21}, AlexaTM 20B for multilingual tasks~\citep{Soltan:22}, OPT~\citep{Zhang:22opt}, Chinchilla~\citep{Hoffmann:22short}, BLOOM~\citep{Scao:22arxivshort}, GLM-130B~\citep{Zeng:22arxivshort}, LaMDA~\citep{Thoppilan:22short}, PaLM~\citep{Chowdhery:22short}, Llama~\citep{Touvron:23arxiv}, and Gemini~\citep{GeminiTeam:23arxiv,Reid:24arxiv}.

\section{Experiments}
\label{sec:Exp}

In this section, we present an empirical analysis of xLSTM, benchmarking it against established approaches in the domain of language modeling. The unique characteristics of xLSTM are explored using synthetic benchmarks in Section~\ref{sec:ExpSyntethic}. Section~\ref{sec:ExpComparison} reports on the validation perplexity for a variety of recent language modeling architectures, each trained on a 15B token subset of the SlimPajama corpus~\citep{Soboleva:23}. We also include ablation experiments to dissect the contributions of different xLSTM components using the same training corpus. Additionally, we investigate scaling properties across all models, following the methodologies of \citet{Kaplan:20} and \citet{Brown:20short}. In Section~\ref{sec:ExpLanguage}, a comprehensive suite of language modeling assessments is conducted. Here, we contrast the performance of xLSTM with that of the top models from Section~\ref{sec:ExpComparison}, all retrained on a 300B token subset from SlimPajama~\citep{Soboleva:23}. The evaluation framework consists of four parts: evaluating extrapolation capabilities to longer sequence contexts, examining validation perplexity alongside downstream application performance as in~\citep{Sutawika:24}, testing generalization on 571 text domains with the PALOMA language benchmark~\citep{Magnusson:23arxivshort}, and analyzing scaling trends across methods, now leveraging a training corpus twenty times larger than in prior comparisons.

Throughout all experiments, we denote xLSTM[$a$:$b$] as indicating an $a/b$ proportion of mLSTM-based to sLSTM-based xLSTM blocks. For instance, xLSTM[7:1] signifies that seven out of eight blocks are mLSTM-based, with the remaining one block being sLSTM-based. In the case where there are 48 blocks in total—a common configuration—this corresponds to 42 mLSTM-based blocks and 6 sLSTM-based blocks. Additionally, in all experimental setups, we incorporate up-projection blocks before mLSTM blocks and after sLSTM blocks.

\subsection{Synthetic Tasks and Long Range Arena}
\label{sec:ExpSyntethic}
We begin by evaluating the impact of xLSTM's exponential gating combined with memory mixing on formal language tasks, following~\citet{Deletang:23}. Next, we examine how well the newly introduced matrix memory in xLSTM performs on the Multi-Query Associative Recall benchmark, as described by~\citet{Arora:23arxiv}. Lastly, we investigate xLSTM's capabilities in handling extended sequences by measuring its results on the Long Range Arena benchmark~\citep{Tay:21}.

\paragraph{Evaluation of xLSTM's Exponential Gating with Memory Mixing.} We assess the performance of xLSTM's novel exponential gating mechanism, enhanced with memory mixing, which is hypothesized to equip the model with the capability to handle state tracking challenges~\citep{Merrill:24, Merrill:23}. To this end, we adopt and modify the formal language benchmarks introduced by \citet{Deletang:23}, adapting them to accommodate training on sequences of various lengths for the purpose of evaluating extrapolation capabilities. Comprehensive explanations of these tasks, along with additional experimental outcomes, are provided in Appendix~\ref{sec:appExpSynthetic-formalLang}. Our study conducts a comparison between xLSTM and several alternative architectures, such as Transformers, State Space Models, and Recurrent Neural Networks. We compute task-specific token accuracy for each method, normalizing results so that scores span from 0 (equivalent to random guessing) to 1 (indicating perfect prediction). For these tasks, we evaluate 2-block variants of the following models: xLSTM[0:1] (representing pure sLSTM), xLSTM[1:0] (pure mLSTM), xLSTM[1:1], Llama, Mamba, RWKV, Retention, Hyena, LSTM, and an LSTM within Transformer blocks (LSTM (Block)). Experimental findings are presented in Figure~\ref{fig:formal-main}. Notably, architectures such as Transformers and State Space Models that lack memory mixing and therefore lack state tracking capabilities are unable to resolve even straightforward regular grammar tasks, such as the parity problem. This aligns with prior results, which indicate a fundamental expressive limitation of Transformers and State Space Models relative to Recurrent Neural Networks~\citep{Merrill:24, Merrill:23, Deletang:23}.

\paragraph{Test of xLSTM's Memory Capacities on Associative Recall Tasks.} This study evaluates the matrix-based memory mechanism of xLSTM by assessing its capacity on the Multi-Query Associative Recall benchmark~\citep{Arora:23arxiv}. In each test sequence, a set of key-value pairs, drawn at random from an extensive vocabulary, must be stored and accurately retrieved later in the sequence. To make the task more challenging than its standard formulation, we scale the number of key-value pairs up to 256 and increase the context length as far as 2048 tokens. This enables a more comprehensive assessment of various models' memory capacities. We conduct comparisons across several 2-block models, namely Llama, Mamba, RWKV-5, RWKV-6, xLSTM[1:1], and xLSTM[1:0]. Model performance is gauged by the proportion of correct pair recalls. Given that Transformer architectures (such as Llama) possess memory capabilities that scale exponentially with the model's representation dimension~\citep{Ramsauer:21}, they serve as the performance upper bound on this task. Outcomes, depicted in Figure~\ref{fig:mqar-main}, show that xLSTM[1:1] outperforms all non-Transformer baselines, maintaining this advantage even for compact model variants. Notably, the presence of the sLSTM block does not reduce the overall memory capability, but actually appears to enhance it, especially evident in the most complex condition with 256 key-value pairs. Further experimental findings are discussed in Appendix~\ref{sec:appExpSynthetic-mqar}, where extrapolation studies suggest that the extended memory faculties of xLSTM facilitate generalization to sequence lengths that surpass those encountered in training.

\paragraph{Test of xLSTM's Long Context Capabilities on Long Range Arena.} To evaluate how xLSTM manages extended sequences and substantial context lengths, we conduct a comparative analysis of several approaches using the Long Range Arena~\citep{Tay:21}. xLSTM consistently achieves robust results across all evaluated tasks, indicating that its architecture effectively addresses various challenges associated with processing long contexts.

Additional information can be found in Appendix~\ref{sec:appExpSynthetic-lra}.

\subsection{Method Comparison and Ablation Study}
\label{sec:ExpComparison}

In order to investigate the central research question—specifically, the performance of our novel LSTM variants when scaled for language modeling—we train xLSTMs, Transformers, State Space Models, along with additional approaches using 15B tokens from SlimPajama within a consistent auto-regressive framework. We evaluate these models on the validation set and conduct ablation experiments focused on the xLSTMs.

\paragraph{Comparing xLSTM to Other Methods.} For our comparison, we train each model on a dataset comprising 15B tokens from SlimPajama~\citep{Soboleva:23}. Model evaluation is based on perplexity calculated over the validation split. The comparison encompasses several approaches: our proposed xLSTM, GPT-3 (Transformer)~\citep{Brown:20short}, Llama (Transformer)~\citep{Touvron:23arxiv}, H3 (SSM)~\citep{Fu:23}, Mamba (SSM)~\citep{Gu:24arxiv}, RWKV-4 (RNN)~\citep{Peng:23arxivshort}, RWKV-5 (RNN)~\citep{Peng:24arxivshort}, RWKV-6 (RNN)~\citep{Peng:24arxivshort}, GLA (linear Transformer)~\citep{Yang:23arxiv}, HGRN (RNN)~\citep{Qin:23}, HGRN2 (RNN)~\citep{Qin:24arxiv}, RetNet (linear Transformer)~\citep{Sun:23arxiv}, Hyena (linear Transformer)~\citep{Poli:23}, xLSTM[1:0], and xLSTM[7:1] (see Section~\ref{sec:xlstm_blocks_notation}). Mixed-precision training was applied across the board; however, for RWKV-5, RWKV-6, GLA, and HGRN2, training did not employ PyTorch's automatic mixed-precision (for more on this, refer to Appendix Section~\ref{sec:appGenTrainProc}).

The evaluated architectures are grouped as follows: (a) Transformers, (b) State Space Models (SSMs), and (c) Recurrent Neural Networks (RNNs) alongside linear Transformers, where the latter category includes models employing linearized attention mechanisms. All architectures are parameter-matched to a 350M-parameter GPT-3 configuration, using an embedding dimension of 1024 and 24 residual blocks. Of note, GPT-3 is unique in using tied weights for token and output embeddings, which slightly reduces its parameter count relative to others. 

Results, summarized in Table~\ref{tab:spaj15b_model_results}, demonstrate that xLSTM achieves the lowest validation perplexity among all compared methods. For a detailed breakdown, consult Appendix~\ref{sec:appExpComparison}. The scaling trends for this experiment are depicted in Figure~\ref{fig:temp_sclaw15B}, suggesting that xLSTM maintains its advantage even as model size increases.

\paragraph{Ablation Studies.}
Table~\ref{tab:spaj15b_model_results} and Figure~\ref{fig:temp_sclaw15B} illustrate that xLSTM attains strong language modeling performance when trained on 15B tokens drawn from the SlimPajama dataset. To systematically dissect the enhancements from LSTM to xLSTM, we incrementally transform a baseline LSTM architecture into the xLSTM framework. The initial step involves embedding LSTM layers within pre-LayerNorm residual structures. In the following phase, we augment the architecture with a post up-projection module. The final modification incorporates exponential gating mechanisms along with a matrix memory component.

Table~\ref{tab:ablstudies} (top) presents these findings.

Our ablation studies indicate that the notable gains in performance arise from the combined contributions of exponential gating and matrix memory. Given the critical role of gating within RNNs and State Space Models, we further evaluate alternative gating strategies. The outcomes, detailed in Table~\ref{tab:ablstudies} (bottom), reveal that when each gate is learnable and conditioned on the input, there is a stepwise enhancement in effectiveness. Further experiments examining the impact of specific backbone components can be found in Appendix~\ref{sec:appAblDetails}.

\subsection{xLSTM as Large Language Model}
\label{sec:ExpLanguage}

This investigation concludes with extensive language modeling experiments aimed at evaluating xLSTM's viability as a large language model (LLM). To this end, we scale up the training dataset, utilizing 300B tokens drawn from SlimPajama. This token count aligns with that employed in prior studies, such as Mamba~\citep{Gu:24arxiv} and Griffin~\citep{De:24arxiv}.

We evaluate xLSTM in relation to RWKV-4, Llama, and Mamba, choosing a representative from each method class discussed in Section~\ref{sec:ExpComparison}. RWKV-4 serves as the RNN exemplar because suitable training precision configurations for RWKV-5, RWKV-6, and HGRN2 were established only after initiating the 300B token experiments (refer to Appendix~\ref{sec:appGenTrainProc} for details). Our experiments encompass multiple model scales (125M, 350M, 760M, and 1.3B parameters), and we analyze each model's ability to extrapolate to longer sequences, as well as their validation performance. Further, we benchmark these models on downstream tasks, evaluate their proficiency in language modeling across 571 distinct text domains from the PALOMA benchmark, and systematically study their scaling law characteristics.

\paragraph{Sequence Length Extrapolation.}
\label{sec:spaj300B_ctx_extrapolation}
We begin by evaluating the ability of large-scale 1.3B-parameter models—specifically, xLSTM, RWKV-4, Llama, and Mamba—to generalize to longer sequence lengths. These models are trained with a context window of 2048 tokens, and subsequently assessed at extended context lengths, reaching up to 16384 tokens. The outcomes of these experiments are illustrated in Figure~\ref{fig:spaj300B_seq_extrapolation}. Notably, among the evaluated approaches, xLSTM consistently sustains low perplexity scores when processing substantially longer contexts.

\paragraph{Validation Perplexity and Downstream Tasks.}
\label{sec:spaj_downstream_eval}
Additionally, across each model scale, we assess xLSTM, RWKV-4, Llama, and Mamba by their next-token prediction accuracy on the SlimPajama validation split, as well as their effectiveness on downstream benchmarks designed to evaluate common sense reasoning capabilities.

The third column in Table~\ref{tab:lmeval} presents the perplexity scores on the validation set for each evaluated method. Notably, both xLSTM[1:0] and xLSTM[7:1] achieve the lowest validation perplexity across all model scales. The subsequent columns of Table~\ref{tab:lmeval} show results on a range of downstream benchmarks. Across nearly every task and for all parameter regimes, xLSTM consistently outperforms alternative approaches—with the sole exception being the ARC benchmark, where Mamba occasionally attains superior performance. Further information can be found in Appendix~\ref{sec:appExpLanguage}.

\paragraph{Performance on PALOMA Language Tasks.} In the third part of our evaluation, we assess the next-token prediction capabilities of xLSTM, RWKV-4, Llama, and Mamba models—each with varying parameter counts—on the PALOMA suite of language tasks~\citep{Magnusson:23arxivshort}. For this, we compute the perplexity scores across 571 distinct text domains, which span sources such as nytimes.com to Reddit's r/depression community. Table~\ref{tab:ppleval} organizes perplexity results by traditional language modeling tasks (first seven columns) and by domain-specific challenges (final five columns). Among the xLSTM variants, xLSTM[1:0] consistently yields lower perplexity compared to xLSTM[7:1] for these tasks.

xLSTM[1:0] demonstrates superior perplexity scores across 568 out of 571 text domains (99.5\%) when compared with Mamba, surpasses Llama in 486 of 571 domains (85.1\%), and outperforms RWKV-4 in 570 out of 571 domains (99.8\%). Further information can be found in Appendix~\ref{sec:appExpLanguage}.

\paragraph{Scaling Laws.} Fourth, we investigate the manifestation of power-law scaling, which enables prediction of model performance at greater parameter scales~\citep{Kaplan:20,Brown:20short}. Figure~\ref{fig:temp_sclaw300B} illustrates this scaling relationship. While all architectures display broadly comparable scaling trends, their performance levels differ by a fixed offset. Among them, RWKV-4 shows the poorest results, trailed by Llama and then Mamba. xLSTM surpasses Mamba, with a performance gap analogous to that observed between Mamba and Llama. These findings suggest that as model size increases, xLSTM is likely to retain its performance edge relative to both Transformer and State-Space architectures.

\textbf{Generation Times and Maximal Throughput.} In our evaluation, depicted in Figure~\ref{fig:inference_time_generation}, we assess the latency of text generation, as well as the maximal throughput as shown in Figure~\ref{fig:inference_time_decoding_throughput} (left), for the various xLSTM models with a parameter count of 1.3B. We benchmark these results against equivalently sized implementations of Mamba, Llama, and RWKV available from HuggingFace, including an implementation of Llama featuring a static key-value cache. During our experiments, neither full cache compilation for the Transformer architecture nor successful compilation of Mamba using \texttt{torch.compile} were feasible. All models were evaluated with a batch size of 1 and a pre-fill of 16 tokens in the text generation setting, an arrangement that should provide the most advantage to the Transformer-based Llama model.

As illustrated in Figure~\ref{fig:inference_time_generation}, xLSTM, along with other recurrent models such as Mamba and RWKV-4, exhibits linear time complexity for text generation, in contrast to the quadratic scaling observed for Llama. For the maximal decoding throughput experiments, we examine varying batch sizes and pre-fill lengths on the Llama model.

In Figure~\ref{fig:inference_time_decoding_throughput} (right), we observe that xLSTM supports substantially larger batch sizes than Llama, a benefit resulting from xLSTM’s constant memory requirements; consequently, xLSTM achieves the greatest throughput across all evaluated configurations.

\section{Limitations}

(i) Unlike mLSTM, the sequential mixing mechanism employed by sLSTM restricts the ability to carry out operations in parallel, thereby precluding the possibility of rapid parallel execution. However, we have implemented an efficient CUDA kernel for sLSTM, which currently achieves speeds less than twice as slow as our parallelized mLSTM kernel. (ii) The present mLSTM CUDA kernels have not undergone extensive optimization, which results in our current setup running approximately four times slower than frameworks such as FlashAttention or the scanning procedure implemented in Mamba. Optimized CUDA kernels, similar to those found in FlashAttention, could significantly improve performance. (iii) mLSTM's matrix memory results in substantial computational complexity, as it involves manipulating $d \times d$ matrices. Despite this, the processes for memory updating and access remain parameter-free and can be efficiently parallelized using standard matrix operations, leading to only a slight increase in total computation time attributable to memory complexity. (iv) Careful consideration must be given to the initialization strategy of the forget gates. (v) Because the matrix memory in mLSTM is decoupled from sequence length, expanding the context window may risk overloading the memory when handling much longer sequences; nonetheless, this issue has not manifested for sequence lengths up to 16k tokens, as detailed in Section~\ref{sec:spaj300B_ctx_extrapolation}. (vi) Given the substantial computational requirements associated with large-scale language model experiments, neither the network architecture nor the hyperparameters—especially for larger xLSTM models—have been comprehensively optimized at this stage. Achieving optimal performance from xLSTM likely requires further dedicated optimization efforts.

\section{Conclusion}
We have provided a partial response to the fundamental question: To what extent can LSTM be scaled for language modeling when expanded to billions of parameters? The evidence thus far allows us to state: ``At least to the level achieved by contemporary architectures such as Transformers or State Space Models.'' By introducing exponential gating with memory mixing and a novel memory framework, we advanced LSTM to xLSTM. Our experiments demonstrate that xLSTM achieves competitive, and in some cases superior, results in language modeling relative to leading approaches such as Transformers and State Space Models. Observed scaling laws suggest that scaling xLSTM further could enable it to rival existing Large Language Models that rely on the Transformer architecture. Beyond language modeling, xLSTM also holds significant promise for other domains, including Reinforcement Learning, Time Series Prediction, and the simulation of physical systems.

\end{document}